%%%%%%%%%%%%%%%%%%%%%%%%%%%%%%%%%%%%

% \subsection{Provenance Annotations}
% We annotate tuples with expressions that are used to condition whether the tuple appears in the ranking resulting from a query. These expressions are made by combining {\it provenance variables}. 
% \begin{definition}[Provenance variables \cite{MLJ22}]
%     Given an associated attribute $A$ from $\attr(D)$ and an associated value $v$ from the domain of the attribute $\domain(\pi_A(D))$ in a database $D$, a {\it provenance variable} is a variable of the form $A_v$.
% \end{definition}
% \begin{example}
%     For example, a provenance variable generated by the \running{} might be $Ma_{CS}$, where $Ma$ is the attribute `Major' and its associated value is `CS', which is in the domain of `Major'. 
% \end{example}
% In our problem, we concern ourselves only with refining the predicates of attributes that are present in the original query. This makes the provenance variables for attributes not in this set irrelevant for deciding whether a tuple is selected, so we are able to restrict the set of provenance variables to attributes only present in the original query. Therefore, we define the set of provenance variables $\provvars$ with respect to a query $Q$ and a database $D$ as the set $\{ A_v \mid \forall A \in \attr(\conds(Q)), v \in \domain(\pi_A(D)) \}$.
% \begin{example}
%     The set of provenance variables for the \running{} given the database in \Cref{ex:running} is $\{Ma_{CS}, Ma_{EE},\allowbreak Ma_{ME}, Hr_{110}, Hr_{100}, Hr_{95}, Hr_{80}, Hr_{75}, Hr_{60}\}$, abbreviating {\tt Major} as $Ma$ and {\tt Hours} as $Hr$.
% \end{example}
% Notice that a query $Q$ can combine tables together, so its not enough to annotate tuples on the base relation level. Instead, we must look at the query if it were to perform no filtering and annotate its output over the database. If we are refining a query $Q$, we let $\tilde{Q}$ be $Q$ with no predicates (i.e. it matches all tuples) and without {\tt LIMIT}. Then, we annotate each tuple from $\tilde{Q}(D)$ with $\sum_{A \in \attr(\conds(Q))} A_{t.A}$, which we denote $\prov(t)$. 
% \begin{example}
%  \label{ex:provexpr}
%     Let $Q$ be the \running{}. Then, $\tilde{Q}$ is {\tt SELECT * FROM Candidates NATURAL JOIN Internships ORDER BY Gpa DESC}. Now, let $t$ be the tuple of the first candidate shown in \Cref{tab:candidates}, then for the query $Q$ in our running example, $\prov(t) = Ma_{CS} + Hr_{100}$.
% \end{example}
% In the context of a query $Q$, a provenance variable takes the value $1$ if its value is matched by all predicates for the attribute. Formally, $A_v$ is $1$ if and only if for every predicate $p$ in $\conds(Q)$ such that $p.A = A$, then either $v\ p.\diamond\ p.C$ if $p$ is numerical, or $v \in p.C$ if $p$ is categorical, and 0 in any other case. Notice that for a query $Q$ and a tuple $t$, if every provenance variable in $\prov(t)$ is $1$, then it must be in the output ranking as the tuple matches every predicate in $Q$. We construct a conditional table (c-table) from $\tilde{Q}(D)$, where we annotate a tuple $t$ with a local condition $\prov(t) = |Attr(Q)|$, conditioning the existence of $t$ in the output the ranking generated by $Q(D)$.
% \begin{example}
%     Let $\prov(t)$ be the provenance expression in \Cref{ex:provexpr} for $Q$, the \running{}. Then, in the context of $Q$, $Ma_{CS}$ is $1$, as $Q$ has one predicate over {\tt Major}, and $CS$ is included in its set of constants $C$. The provenance variable $Hr_{100}$ is also $1$, as $Q$ has one predicate over {\tt Hours} filtering tuples with ${\tt Hours} \geq 90$, and $100 \geq 90$. Finally, $\prov(t) = 2 = |Attr(Q)|$, implying that the tuple must be included in the output of $Q(D)$.
% \end{example}
% A possible world is the c-table resulting from a particular assignment of the provenance variables in $\provvars$, which are assigned based on the predicates of the (refined) query. 

%\indent \textit{Ranking.} We now develop the variables and constraints needed in order to have a ranking ordered by the {\tt ORDER BY} expression of the input query $Q$. Our model must accommodate a ranking of the maximum size attainable by refining predicates of $Q$. This is ranking is generated by the refinement over $Q$ that includes all tuples, i.e. it effectively has no predicates. Therefore, we let $\tilde{Q}$ be the query $Q$ with its predicates removed. Bearing this in mind, in order to model a ranking, we create a set of variables $x_{i, t} \in \{0, 1\}$ for each position $i$ in $[|\tilde{Q}(D)|]$ and tuple $t$ in $\tilde{Q}(D)$. When $x_{i, t}$ is 1, tuple $t$ has been placed into position $i$. A ranking may not have more than one tuple assigned to a single position, which we model with the constraint $$\sum_{t \in \tilde{Q}(D)} x_{i, t} \leq 1$$ Furthermore, when ranking according to the {\tt ORDER BY} expression of the input query $Q$, we must ensure that each tuple in position $i$ is ranked higher than the tuple in position $i+1$ according to the {\tt ORDER BY} expression of $Q$. In order to do so, we define a set of variables $s_i$ for each position $i$ in the ranking. It is considerably simpler to model its value as $0$ when no tuple is assigned at position $i$. However, modeling this comes with the caveat that we must then consider the tuples in the reverse order of $\tilde{Q}(D)$. Otherwise, we would have to deal with the fact that $0$ would rank higher than anything else in ascending order. Therefore, the value of $s_i$ when there is a tuple present is the rank of the tuple in position $i$ in the reverse order of $\tilde{Q}(D)$. We model this by adding the constraint $$\sum_{t \in \tilde{Q}(D)} (|\tilde{Q}(D)| - \tilde{Q}(D)(t) + 1) \cdot x_{i, t} - s_i = 0$$ for each position $i$ in $[|\tilde{Q}(D)|]$. All that is left is to ensure that the tuples are ordered according to the {\tt ORDER BY} of $Q$. This is the case when every $s_i$ is larger than that of the succeeding position $s_{i+1}$. We model this with the constraint $$s_{i} - s_{i+1} \geq 0$$ for every position $i$ in $[|\tilde{Q}(D)|]$. Effectively, we ensure that the tuples are in descending rank of the reversed $\tilde{Q}(D)$, which is equivalent to ascending rank in the original ranking. These constraints limit feasible solutions to rankings that are ordered according to the {\tt ORDER BY} of $Q$.
%\begin{example}
%Consider assigning just candidates 1 and 5 from the \running{} to some positions in the ranking. If candidate 1 is placed anywhere but the first position, then above it would be either candidate 5 or no tuple. Let candidate 1 be placed at position $i+1$ where $1 \leq i \leq |\tilde{Q}(D)| - 1$, then $s_{i+1}$ would be $8$ by the constraint above. Therefore, $s_i$ would be either $4$ or $0$ by the same constraint. In either case, this violates the constraint that $s_{i} - s_{i+1} \geq 0$. The only satisfying assignment is to therefore place candidate 1 in position 1, i.e. $x_{1, 1} = 1$. Candidate 5 must then be in position 2, as otherwise the position above it would be empty and therefore have a 0 $s_i$ value, violating the same constraint. Therefore, the only feasible assignment of these tuples to positions is to have candidate 1 in position 1 and candidate 5 in position 2, i.e. $x_{1, 1} = 1$ and $x_{2, 5} = 1$.  
%\end{example}

%\indent \textbf{Variables.} Recall that a tuple is selected by a query if each predicate in the query selects the value the tuple has for that attribute. So, we must decide for each value whether it is selected by the current set of refined predicates. We introduce a variable $A_{v} \in \{0, 1\}$ for each attribute with predicates in $Q$ and each value $v$ in the domain of $A$ in the input database. Then, we create a variable $r_t \in \{0, 1\}$ for each tuple $t$ that could appear in any refinement made to $Q$, i.e. every tuple in $\tilde{Q}(D)$. The variable $r_t$ is $1$ when the tuple $t$ is selected, i.e. all the values of the attributes with predicates are selected by the refinement. In order to ensure that a higher ranked tuple is not placed below a lower ranked one, we must know the rank of the tuple placed in position $i$. Therefore, we create a set of variables $s_i$ where its value represents, for position $i$ in the possible world, the position of the tuple in $\tilde{Q}(D)$. If no tuple is assigned to position $i$, the value of $s_i$ is 0. \\
%\indent \textbf{Constraints.} Given these variables, we must now define constraints over them in order to enforce the desired behavior. First, we develop the constraints necessary to have $A_v$ be $1$ if and only if the value $v$ is selected by the predicate(s) over $A$ of the refinement. For categorical values, its $A_v$ variable is exactly equivalent to the value of $C_{A, v}$ variable, \resolve{why} and we therefore treat $A_v$ for categorical values as an alias for $C_{A, v}$. However, numerical values lack such a direct correspondence to the $C_{v, \diamond}$ variables for numerical predicates. The $C_{v, \diamond}$ variables only define the endpoint(s) of a selected interval, and do not include values directly as in the categorical case. Therefore, the variable $A_v$ for any numerical attribute $A$ must be $1$ if and only if they are within the interval formed by the refined predicate(s) over $A$. We must then check that the value is selected by each individual predicate. Towards this end, we define a set of auxiliary variables $A_{v, \diamond}$ for each predicate $A \diamond C \in \num(Q)$, where $A_{v, \diamond}$ if the value $v$ is selected by this particular predicate. In order to model this, we employ a well-known modeling trick where we let $M_A$ be a constant larger than the maximum magnitude value of some attribute $A$. For lower-bound predicates, to enforce $A_{v, \diamond}$ is $1$ when a value $v$ is selected by the individual predicate, we then add two constraints for every value $A_{v,\diamond}$: 
% \begin{align*}
%     C_{A, \diamond} + M_A \cdot (1 - A_{v, \diamond}) \geq v - {\sf Eq}(\diamond) \cdot \delta \\ 
%     C_{A, \diamond} - M_A \cdot A_{v, \diamond} \leq v - {\sf Eq}(\diamond) \cdot \delta
% \end{align*} where ${\sf Eq}(\diamond)$ is $1$ if $\diamond$ is a strict inequality and $0$ otherwise, and $\delta$ is some small constant to relax the constraint as MILP constraints do not support strict inequalities. Upper-bound predicates are treated similarly, where we instead add the constraints
% \begin{align*}
% C_{A, \diamond} + M_A \cdot (A_{v, \diamond}) \geq v + {\sf Eq}(\diamond) \cdot \delta \\
% C_{A, \diamond} - M_A \cdot (1 - A_{v, \diamond}) \leq v + {\sf Eq}(\diamond) \cdot \delta
% \end{align*}
% If every numerical predicate over $A$ selects the value $v$, i.e. every $A_{v, \diamond}$ is 1 for a value $v$, then $v$ is selected by the refined predicates over $A$, and $A_v$ must be $1$. Let $P$ be the set of numeric predicates over $A$, then we model the behavior of $A_v$ for numerical values with the constraint $$0 \leq \sum_{(A \diamond C) \in P} A_{v, \diamond} - |P| \cdot A_{v} < |P|$$
% \begin{example}
% Consider refining the \running{} over the database in \Cref{ex:running}. Let $C_{Hours, \geq}$ be 80. Then, \resolve{example}.
% \end{example}
% \indent Now that have constraints that decide which values are selected by the refined predicates, we must determine which tuples are selected. Recall that a tuple is selected if and only if all its values for each attribute match the predicate(s) over those attributes. Therefore, we annotate each tuple $t$ with a set of $A_v$ variables denoted by $\prov(t)$, which we define as $\{A_{t.A} \mid \forall A \in \attr(\conds(Q)) \}$. Then, for a tuple $t$ we count the number of selected variables in $\prov(t)$. If all are selected, i.e. there are $|\attr(\conds(Q))|$ selected variables, then the tuple is included in the ranking as it is selected by the all the refined predicates. We use this annotation to decide the value of $r_t$ by adding the constraint $$0 \leq \sum_{A_v \in \prov(t)} A_v - |\attr(\conds(Q))| \cdot r_t < |\attr(\conds(Q))|$$ for each tuple.
% \begin{example}
% Let $t$ be the first tuple presented in \Cref{tab:candidates} joined with \Cref{tab:internships} according to the \running{}. Then, its lineage is the set of variables $\{Major_{CS}, Hours_{100}\}$. Consider the refinements made by $C_{Major, CS} = 1$, $C_{Major, EE} = 1$ and $C_{Hours, \geq} = 80$. Then, all the variables in the tuple's lineage would take the value $1$. The variable $Major_{CS} \equiv C_{Major, CS}$ has value $1$ by our supposition. \resolve{The variable} $Hours_{100}$ takes the value $1$ as $100$ is selected by all (the only) predicates over ${\tt Hours}$, in this case ${\tt Hours} \geq 80$. This makes the number of selected values $2$, and as $r_t \in \{0, 1\}$, the only feasible assignment to $r_t$ is $1$, meaning the tuple appears in the ranking as desired.  
% \end{example}
% \indent At this point we have selected the tuples that are selected by the current refinement, and we need to assign them in positions of the ranking. By specifying  for each tuple, we enforce that every tuple selected by the current refinement must be assigned to some position in the ranking. In order to not assign more than one tuple to the same position, we enforce  for each position $i$ in the ranking. These MIP constraints together enforce condition \numbercircled{1}. To enforce condition \numbercircled{2}, for each position $i$, we equate $s_i$ with the position of the tuple placed into position $i$ in $\tilde{Q}(D)$ with the constraint: $$\sum_{t \in \tilde{Q}(D)} (\tilde{Q}(D))(t) \cdot (x_{i, t}) - s_i = 0$$ We now enforce that higher ranked tuples come before lower ranked tuples in our possible world with the constraint $$s_{i+1} - s_i \geq 0$$ for every position up to $i-1$.
% \begin{example}
%     \resolve{example}
% \end{example}
% \begin{lemma}
%     For a query $Q$ and database $D$, the ranking $Q(D)$ is equivalent to the ranking $[t \mid \forall i \in [|\tilde{Q}(D)|], \exists x_{i, t}, x_{i, t} = 1]$ for a feasible assignment of variables to the MILP problem.
% \end{lemma}
% \begin{proof}
%     \resolve{proof}
% \end{proof}

% old opt
% The first takes advantage of the fact that for a constraint set $\constraints$ only the first $k^*$ tuples matter in determining the deviation of the ranking from the constraint set. In the original mixed-integer linear program, we enforce all the possible world conditions over all positions in the ranking. However, we can relax some constraints by only enforcing them over the first $k^*$ tuples in the ranking. Specifically, altering \Cref{eq:score_integrity} to be for all $i$ in $[k^* - 1]$, \Cref{eq:position_integrity} to be for all $i$ in $[k^*]$, \Cref{eq:rank_tuples_in_prov} to sum over $i$ in $[k^*]$, and \Cref{eq:set_rank_score} to be for all $i$ in $[k^*]$. As we do not need to check tuples past the $k^*$ position in order to ensure correctness of the top-$k^*$, we can relax these constraints to make finding solutions easier.
% \begin{theorem}
%     For a query $Q$ and database $D$, the ranking $Q(D)_{k^*}$ is equivalent to the ranking $[t \mid \forall i \in [k^*], \exists x_{i, t}, x_{i, t} = 1]$ given a feasible assignment of variables for the version of the MIP problem with these optimized constraints.
% \end{theorem}
% \begin{proof}
%     \resolve{proof}
% \end{proof}

%\subsubsection{Recovering the Refinement}
%\indent After solving our mixed-integer linear program, we must recover the refinement from the values of the variables in the solution to answer \problem{}. Recall that the values of categorical attributes are selected by predicates that select the values if they are in the predicates' set of constants $C$. Therefore, the values $v$ of the $C_{A, v}$ variables that are $1$ for a categorical attribute $A$ form the set of constants $C$ for the new refined predicate over $A$. In the case of numerical attributes, we refine the constant $C$ of a predicate $A \diamond C$ to the value in $C_{A, \diamond}$. The output refinement is then the predicates refined to the (set of) constants as described here, answering \problem{}.